% answers.tex —— 2016 年·上海卷·文科 · 参考答案与详细解析片段
% 由 tools/generate_projects.py 依来源总账自 2016/2016_数学真题_上海文_试题.md + 2016/2016_数学真题_上海文_解析正文.md + 2016/2016_数学真题_上海文_答案要点.md 生成，勿手改（改源 MD 后重跑）


\begin{gktip}
固定原卷不附官方答案。以下为本库独立推导，非官方标准答案；第11题按可核验公开解析采用古典概型的通常解读，并独立复算。
\end{gktip}

\begin{answer}
\gksolhead{1}{}

\textbf{题目}：

设 \(x\in\mathbf R\)，则不等式 \(|x-3|<1\) 的解集为\(\underline{\qquad}\)。

\Needspace{4\baselineskip}
\textbf{解析}

答案：\((2,4)\)。由 \(|x-3|<1\) 得 \(-1<x-3<1\)，故 \(2<x<4\)。

\gkanswer{答案：\( (2,4) \)}

\end{answer}
\begin{answer}
\gksolhead{2}{}

\textbf{题目}：

设 \(z=\dfrac{3+2\mathrm i}{\mathrm i}\)，其中 \(\mathrm i\) 为虚数单位，则 \(z\) 的虚部等于\(\underline{\qquad}\)。

\Needspace{4\baselineskip}
\textbf{解析}

答案：\(-3\)。\(z=(3+2i)/i=(3+2i)(-i)=2-3i\)，虚部为实数 \(-3\)。

\gkanswer{答案：\(-3\)}

\end{answer}
\begin{answer}
\gksolhead{3}{}

\textbf{题目}：

已知平行直线 \(l_1:\ 2x+y-1=0\)，\(l_2:\ 2x+y+1=0\)，则 \(l_1\) 与 \(l_2\) 的距离是\(\underline{\qquad}\)。

\Needspace{4\baselineskip}
\textbf{解析}

答案：\(2/\sqrt5\)。两平行直线法向量为 \((2,1)\)，常数项相差 \(2\)，距离为 \(|1-(-1)|/\sqrt{2^2+1^2}\)。

\gkanswer{答案：\(\dfrac{2}{\sqrt5}\)}

\end{answer}
\begin{answer}
\gksolhead{4}{}

\textbf{题目}：

某次体检，5 位同学的身高（单位：米）分别为 1.72，1.78，1.80，1.69，1.76，则这组数据的中位数是\(\underline{\qquad}\)（米）。

\Needspace{4\baselineskip}
\textbf{解析}

答案：\(1.76\) 米。排序为 \(1.69,1.72,1.76,1.78,1.80\)，中间第三个数即中位数。

\gkanswer{答案：\(1.76\) 米}

\end{answer}
\begin{answer}
\gksolhead{5}{}

\textbf{题目}：

若函数 \(f(x)=4\sin x+a\cos x\) 的最大值为 5，则常数 \(a=\)\(\underline{\qquad}\)。

\Needspace{4\baselineskip}
\textbf{解析}

答案：\(a=\pm3\)。振幅为 \(\sqrt{16+a^2}\)，最大值为5，故 \(a^2=9\)；两种符号均可达到最大值5。

\gkanswer{答案：\(a=\pm3\)}

\end{answer}
\begin{answer}
\gksolhead{6}{}

\textbf{题目}：

已知点 \((3,9)\) 在函数 \(f(x)=1+a^x\) 的图象上，则 \(f(x)\) 的反函数 \(f^{-1}(x)=\)\(\underline{\qquad}\)。

\Needspace{4\baselineskip}
\textbf{解析}

答案：\(f^{-1}(x)=\log_2(x-1)\)，\(x>1\)。代入 \((3,9)\) 得 \(a^3=8\)，所以 \(a=2\)；从 \(y=1+2^x\) 解出 \(x=\log_2(y-1)\)。

\gkanswer{答案：\(f^{-1}(x)=\log_2(x-1)\)，\(x>1\)}

\end{answer}
\begin{answer}
\gksolhead{7}{}

\textbf{题目}：

\Needspace{5\baselineskip}
若 \(x,y\) 满足
\nopagebreak[4]
\[
\begin{cases}
x\geqslant0,\\
y\geqslant0,\\
y\geqslant x+1,
\end{cases}
\]
则 \(x-2y\) 的最大值为\(\underline{\qquad}\)。

\Needspace{4\baselineskip}
\textbf{解析}

答案：\(-2\)。由 \(y\ge x+1\) 得 \(x-2y\le-x-2\le-2\)。在 \((x,y)=(0,1)\) 取等。

\gkanswer{答案：\(-2\)}

\end{answer}
\begin{answer}
\gksolhead{8}{}

\textbf{题目}：

方程 \(3\sin x=1+\cos2x\) 在区间 \([0,2\pi]\) 上的解为\(\underline{\qquad}\)。

\Needspace{4\baselineskip}
\textbf{解析}

答案：\(\{\pi/6,5\pi/6\}\)。令 \(s=\sin x\)，由 \(1+\cos2x=2-2s^2\) 得 \(2s^2+3s-2=(2s-1)(s+2)=0\)，所以 \(s=1/2\) 或 \(-2\)。后者超出正弦值域，前者在 \([0,2\pi]\) 中给出上述两解。
核验：\(\sin x=1/2\) 时左边 \(3/2\)，右边 \(1+\cos2x=3/2\)。

\gkanswer{答案：\(\dfrac{\pi}{6},\dfrac{5\pi}{6}\)}

\end{answer}
\begin{answer}
\gksolhead{9}{}

\textbf{题目}：

在 \(\left(\sqrt[3]{x}-\dfrac2x\right)^n\) 的二项展开式中，所有项的二项式系数之和为 256，则常数项等于\(\underline{\qquad}\)。

\Needspace{4\baselineskip}
\textbf{解析}

答案：\(112\)。二项式系数和为 \(2^n=256\)，所以 \(n=8\)。第 \(k+1\) 项为 \(\binom8k(-2)^k x^{(8-k)/3-k}\)；常数项须 \(8-4k=0\)，得 \(k=2\)，常数为 \(\binom82\cdot4=112\)。

\gkanswer{答案：\(112\)}

\end{answer}
\begin{answer}
\gksolhead{10}{}

\textbf{题目}：

已知 \(\triangle ABC\) 的三边长分别为 3，5，7，则该三角形的外接圆半径等于\(\underline{\qquad}\)。

\Needspace{4\baselineskip}
\textbf{解析}

答案：\(7\sqrt3/3\)。半周长 \(s=15/2\)，海伦公式给面积
\(\Delta=\sqrt{(15/2)(9/2)(5/2)(1/2)}=15\sqrt3/4\)。
因此外接圆半径 \(R=abc/(4\Delta)=105/(15\sqrt3)=7/\sqrt3\)。
数值复算：面积约 \(6.49519\)，半径约 \(4.04145\)，\(4R\Delta=105\)。

\gkanswer{答案：\(\dfrac{7\sqrt3}{3}\)}

\end{answer}
\begin{answer}
\gksolhead{11}{}

\textbf{题目}：

某食堂规定，每份午餐可以在四种水果中任选两种，则甲、乙两同学各自所选的两种水果相同的概率为\(\underline{\qquad}\)。

\Needspace{4\baselineskip}
\textbf{解析}

答案：\(\dfrac16\)。按公开解析采用的古典概型，将甲、乙的选择视为独立，且每人从四种水果的 \(\binom42=6\) 种组合中等可能选择。共有 \(6^2=36\) 个有序结果，其中两人选法相同的有6个，故所求概率为 \(6/36=1/6\)。该模型是公开解析所采用的通常解读；固定原卷不附官方答案。

\gkanswer{答案：\(\dfrac16\)}

\end{answer}
\begin{answer}
\gksolhead{12}{}

\textbf{题目}：

如图，已知点 \(O(0,0)\)，\(A(1,0)\)，\(B(0,-1)\)，\(P\) 是曲线 \(y=\sqrt{1-x^2}\) 上一个动点，则 \(\overrightarrow{OP}\cdot\overrightarrow{BA}\) 的取值范围是\(\underline{\qquad}\)。
\nopagebreak[4]
\begin{center}\begin{tikzpicture}[scale=1.25,>=Stealth,every node/.style={font=\small,inner sep=2pt}]
\coordinate (O) at (0,0); \coordinate (A) at (1,0); \coordinate (B) at (0,-1); \coordinate (P) at (135:1);
\draw[->] (-1.4,0)--(1.45,0) node[right] {$x$};
\draw[->] (0,-1.45)--(0,1.45) node[above] {$y$};
\draw (1,0) arc[start angle=0,end angle=180,radius=1];
\draw[->] (O)--(P); \draw[->] (B)--(A);
\node[below left] at (O) {$O$}; \node[below right] at (A) {$A$}; \node[below left] at (B) {$B$}; \node[above left] at (P) {$P$};
\end{tikzpicture}
\end{center}
\nopagebreak[4]
\Needspace{4\baselineskip}
\textbf{解析}

答案：\([-1,\sqrt2]\)。\(\overrightarrow{BA}=(1,1)\)，令 \(P=(\cos\theta,\sin\theta)\)，\(0\le\theta\le\pi\)，则点积为 \(\cos\theta+\sin\theta=\sqrt2\sin(\theta+\pi/4)\)。最大值在 \(\theta=\pi/4\) 为 \(\sqrt2\)，最小值在 \(\theta=\pi\) 为 \(-1\)；连续性保证中间值全部取得。

\gkanswer{答案：\([-1,\sqrt2]\)}

\end{answer}
\begin{answer}
\gksolhead{13}{}

\textbf{题目}：

\Needspace{5\baselineskip}
设 \(a>0\)，\(b>0\)，若关于 \(x,y\) 的方程组
\nopagebreak[4]
\[
\begin{cases}
ax+y=1,\\
x+by=1
\end{cases}
\]
无解，则 \(a+b\) 的取值范围是\(\underline{\qquad}\)。

\Needspace{4\baselineskip}
\textbf{解析}

答案：\((2,+\infty)\)。无解必须 \(ab-1=0\)。此时第二个方程乘 \(a\) 后为 \(ax+y=a\)，与第一个 \(ax+y=1\) 矛盾当且仅当 \(a\ne1\)。故 \(b=1/a\) 且 \(a>0,a\ne1\)，由 \(a+1/a>2\) 得所求范围。反之每个 \(t>2\) 均可由 \(a^2-ta+1=0\) 的正根实现，故没有遗漏。

\gkanswer{答案：\((2,+\infty)\)}

\end{answer}
\begin{answer}
\gksolhead{14}{}

\textbf{题目}：

无穷数列 \(\{a_n\}\) 由 \(k\) 个不同的数组成，\(S_n\) 为 \(\{a_n\}\) 的前 \(n\) 项和。若对任意的 \(n\in\mathbf N^*\)，\(S_n\in\{2,3\}\)，则 \(k\) 的最大值为\(\underline{\qquad}\)。

\Needspace{4\baselineskip}
\textbf{解析}

答案：\(k_{\max}=4\)。\(a_1=S_1\) 只能为2或3中的一个，\(n\ge2\) 时 \(a_n=S_n-S_{n-1}\in\{-1,0,1\}\)，所以最多四种数。取前六个部分和为 \(2,2,3,2,3,3\)，之后恒为3，则数列开头为 \(2,0,1,-1,1,0\)，其后恒0，恰好出现 \(2,-1,0,1\) 四种，达到上界。

\gkanswer{答案：\(4\)}

\end{answer}
\begin{answer}
\gksolhead{15}{}

\textbf{题目}：

设 \(a\in\mathbf R\)，则“\(a>1\)”是“\(a^2>1\)”的（　）

\noindent(A) 充分非必要条件

\noindent(B) 必要非充分条件

\noindent(C) 充要条件

\noindent(D) 既非充分也非必要条件

\Needspace{4\baselineskip}
\textbf{解析}

答案：A，充分非必要条件。\(a>1\) 必然 \(a^2>1\)；反向还可有 \(a<-1\)，例如 \(a=-2\)。

\gkanswer{答案：A}

\end{answer}
\begin{answer}
\gksolhead{16}{}

\textbf{题目}：

如图，在正方体 \(ABCD-A_1B_1C_1D_1\) 中，\(E,F\) 分别为 \(BC,BB_1\) 的中点，则下列直线中与直线 \(EF\) 相交的是（　）
\nopagebreak[4]
\begin{center}\begin{tikzpicture}[scale=1.05,>=Stealth,every node/.style={font=\small,inner sep=2pt}]
\coordinate (A) at (0,0);\coordinate (B) at (2,0);\coordinate (C) at (2.9,.9);\coordinate (D) at (.9,.9);
\coordinate (A1) at (0,2);\coordinate (B1) at (2,2);\coordinate (C1) at (2.9,2.9);\coordinate (D1) at (.9,2.9);
\coordinate (E) at (2.45,.45);\coordinate (F) at (2,1);
\draw[dashed] (A)--(D)--(C) (D)--(D1);
\draw (A)--(B)--(C)--(C1)--(D1)--(A1)--(A) (A1)--(B1)--(B)--(C) (B1)--(C1);
\fill (E) circle (1pt) (F) circle (1pt);
\node[below left] at (A) {$A$};\node[below] at (B) {$B$};\node[right] at (C) {$C$};\node[below right,xshift=2pt,yshift=-2pt] at (D) {$D$};
\node[left] at (A1) {$A_1$};\node[above left] at (B1) {$B_1$};\node[above right] at (C1) {$C_1$};\node[above] at (D1) {$D_1$};
\node[below right] at (E) {$E$};\node[above right,xshift=2pt,yshift=2pt] at (F) {$F$};
\end{tikzpicture}
\end{center}
\nopagebreak[4]
\noindent(A) 直线 \(AA_1\)

\noindent(B) 直线 \(A_1B_1\)

\noindent(C) 直线 \(A_1D_1\)

\noindent(D) 直线 \(B_1C_1\)

\Needspace{4\baselineskip}
\textbf{解析}

答案：D，直线 \(B_1C_1\)。取 \(A=(0,0,0), B=(1,0,0), C=(1,1,0), A_1=(0,0,1)\)，则 \(E=(1,1/2,0),F=(1,0,1/2)\)。直线 \(EF\) 满足 \(x=1,y+z=1/2\)；直线 \(B_1C_1\) 满足 \(x=1,z=1\)，在 \((1,-1/2,1)\) 相交。交点在边的延长线上，题面问的是直线，故有效。其他三条分别位于 \(x=0\) 或 \(y=0,z=1\)，均不相交。

\gkanswer{答案：D}

\end{answer}
\begin{answer}
\gksolhead{17}{}

\textbf{题目}：

设 \(a\in\mathbf R\)，\(b\in[0,2\pi)\)。若对任意实数 \(x\) 都有 \(\sin\left(3x-\dfrac\pi3\right)=\sin(ax+b)\)，则满足条件的有序实数对 \((a,b)\) 的对数为（　）

\noindent(A) 1

\noindent(B) 2

\noindent(C) 3

\noindent(D) 4

\Needspace{4\baselineskip}
\textbf{解析}

答案：B，共2对：\((3,5\pi/3)\)、\((-3,4\pi/3)\)。共同的非零正弦函数满足二阶微分方程 \(u''=-9u=-a^2u\)，故 \(a=\pm3\)。\(a=3\) 时相位须 \(b\equiv-\pi/3\)；\(a=-3\) 时用 \(\sin v=\sin(\pi-v)\) 得 \(b\equiv4\pi/3\)。在 \([0,2\pi)\) 中各一解。

\gkanswer{答案：B}

\end{answer}
\begin{answer}
\gksolhead{18}{}

\textbf{题目}：

设 \(f(x)\)，\(g(x)\)，\(h(x)\) 是定义域为 \(\mathbf R\) 的三个函数，对于命题：① 若 \(f(x)+g(x)\)、\(f(x)+h(x)\)、\(g(x)+h(x)\) 均为增函数，则 \(f(x)\)、\(g(x)\)、\(h(x)\) 中至少有一个为增函数；② 若 \(f(x)+g(x)\)、\(f(x)+h(x)\)、\(g(x)+h(x)\) 均是以 \(T\) 为周期的函数，则 \(f(x)\)、\(g(x)\)、\(h(x)\) 均是以 \(T\) 为周期的函数，下列判断正确的是（　）

\noindent(A) ①和②均为真命题

\noindent(B) ①和②均为假命题

\noindent(C) ①为真命题，②为假命题

\noindent(D) ①为假命题，②为真命题

\Needspace{4\baselineskip}
\textbf{解析}

答案：D，①假，②真。

①反例：取 \(f=x+2\sin x\)，\(g=x+2\sin(x+2\pi/3)\)，\(h=x+2\sin(x+4\pi/3)\)。每个函数的导数在相位等于 \(\pi\) 附近为负，故都不是增函数；同时各自又有正导数区间，亦不是单调函数。任意两个的和形如 \(2x+2\sin(x+\phi)\)（可能吸收符号到相位），导数为 \(2+2\cos(x+\phi)\ge0\)，且在任意非退化区间不恒为0，故严格递增。①被否定。

②由 \(2f=(f+g)+(f+h)-(g+h)\)，三个和函数均满足 \(u(x+T)=u(x)\)，所以 \(f(x+T)=f(x)\)；同理 \(g,h\)。这里“以T为周期”指T为一个周期，题面未称最小正周期。

\gkanswer{答案：D}

\end{answer}
\begin{answer}
\gksolhead{19}{}

\textbf{题目}：

将边长为 1 的正方形 \(AA_1O_1O\)（及其内部）绕 \(OO_1\) 旋转一周形成圆柱，如图，\(\overset{\frown}{AC}\) 长为 \(\dfrac{5\pi}6\)，\(\overset{\frown}{A_1B_1}\) 长为 \(\dfrac\pi3\)，其中 \(B_1\) 与 \(C\) 在平面 \(AA_1O_1O\) 的同侧。

（1） 求圆柱的体积与侧面积；

（2） 求异面直线 \(O_1B_1\) 与 \(OC\) 所成的角的大小。
\nopagebreak[4]
\begin{center}\begin{tikzpicture}[scale=1,>=Stealth,every node/.style={font=\small,inner sep=2pt}]
\coordinate (O) at (0,0); \coordinate (O1) at (0,2.2);
\coordinate (A) at (1.7,0); \coordinate (A1) at (1.7,2.2);
\coordinate (B1) at (.85,1.72);
\coordinate (C) at (-1.47,-.28);
\draw[dashed] (1.7,0) arc[start angle=0,end angle=180,x radius=1.7,y radius=.56];
\draw (-1.7,0) arc[start angle=180,end angle=360,x radius=1.7,y radius=.56];
\draw (O1) ellipse[x radius=1.7,y radius=.56];
\draw (-1.7,0)--(-1.7,2.2) (A)--(A1) (O1)--(A1) (O1)--(B1);
\draw[dashed] (O1)--(O)--(A) (O)--(C);

\node[below right,xshift=2pt] at (O) {$O$};\node[above left] at (O1) {$O_1$};
\node[right] at (A) {$A$};\node[right] at (A1) {$A_1$};\node[below right] at (B1) {$B_1$};\node[below left] at (C) {$C$};
\end{tikzpicture}
\end{center}
\nopagebreak[4]
\Needspace{4\baselineskip}
\textbf{解析}

\gksubpart{(1)}

答案：体积 \(\pi\)，侧面积 \(2\pi\)。旋转轴为正方形边 \(OO_1\)，半径、高均为1；\(V=\pi r^2h\)，\(S_{\text{侧}}=2\pi rh\)。

\gksubpart{(2)}

答案：\(\pi/2\)。半径1使弧长等于圆心角，\(\angle AOC=5\pi/6\)，\(\angle A_1O_1B_1=\pi/3\)。\(B_1,C\) 在平面 \(AA_1O_1O\) 同侧，取该侧为正转向，则两径向向量的转角分别为 \(\pi/3\) 与 \(5\pi/6\)，差为 \(\pi/2\)。\(O_1B_1\) 与 \(OC\) 均为水平径向方向，内积为0，异面直线所成角为直角。

\gkanswer{答案（1）：体积 \(\pi\)，侧面积 \(2\pi\)。}

\gkanswer{答案（2）：\(\dfrac{\pi}{2}\)。}

\end{answer}
\begin{answer}
\gksolhead{20}{}

\textbf{题目}：

有一块正方形菜地 \(EFGH\)，\(EH\) 所在直线是一条小河，收获的蔬菜可送到 \(F\) 点或河边运走。于是，菜地分为两个区域 \(S_1\) 和 \(S_2\)，其中 \(S_1\) 中的蔬菜运到河边较近，\(S_2\) 中的蔬菜运到 \(F\) 点较近，而菜地内 \(S_1\) 和 \(S_2\) 的分界线 \(C\) 上的点到河边与到 \(F\) 点的距离相等，现建立平面直角坐标系，其中原点 \(O\) 为 \(EF\) 的中点，点 \(F\) 的坐标为 \((1,0)\)，如图。

（1） 求菜地内的分界线 \(C\) 的方程。

（2） 菜农从蔬菜运量估计出 \(S_1\) 面积是 \(S_2\) 面积的两倍，由此得到 \(S_1\) 面积的“经验值”为 \(\dfrac83\)。设 \(M\) 是 \(C\) 上纵坐标为 1 的点，请计算以 \(EH\) 为一边，另一边过点 \(M\) 的矩形的面积，及五边形 \(EOMGH\) 的面积，并判断哪一个更接近于 \(S_1\) 面积的经验值。
\nopagebreak[4]
\begin{center}\begin{tikzpicture}[scale=1.1,>=Stealth,every node/.style={font=\small,inner sep=2pt}]
\draw[->] (-1.6,0)--(1.6,0) node[right] {$x$};
\draw[->] (0,-.55)--(0,2.6) node[above] {$y$};
\draw (-1,0)--(-1,2)--(1,2)--(1,0);
\draw[domain=0:2,samples=80,smooth,variable=\t] plot ({\t*\t/4},{\t});
\fill (.25,1) circle (1pt) node[right] {$M$};
\node[below left] at (0,0) {$O$}; \node[below] at (-1,0) {$E$};\node[below] at (1,0) {$F$};
\node[above left] at (-1,2) {$H$};\node[above right] at (1,2) {$G$};
\node at (-.55,1.25) {$S_1$};\node at (.65,.65) {$S_2$};
\end{tikzpicture}
\end{center}
\nopagebreak[4]
\Needspace{4\baselineskip}
\textbf{解析}

\gksubpart{(1)}

答案：\(y^2=4x\)，菜地内 \(0\le y\le2\)（等价 \(0\le x\le1\)）。\(E=(-1,0)\)，河流为直线 \(x=-1\)。菜地内点到河流距离为 \(x+1\)，到 \(F=(1,0)\) 的距离为 \(\sqrt{(x-1)^2+y^2}\)。平方等式得 \(y^2=4x\)。正方形边长2，所以其内有效段从 \(O\) 到 \(G=(1,2)\)。

\gksubpart{(2)}

\(M=(1/4,1)\)。以 \(EH\) 为一边且另一边过M的矩形，宽 \(1+1/4=5/4\)，高2，面积 \(5/2\)。
五边形 \(EOMGH\) 的面积可沿高度分段求：折线 \(OM\) 的横坐标为 \(x=y/4\)（\(0\le y\le1\)），\(MG\) 为 \(x=3y/4-1/2\)（\(1\le y\le2\)）。故面积
\(\int_0^1(1+y/4)dy+\int_1^2(1+3y/4-1/2)dy=9/8+13/8=11/4\)。
与经验值 \(8/3\) 的差分别为 \(1/6\) 与 \(1/12\)，所以 \textbf{五边形更接近}。数值为 \(2.5,2.75,2.666666\ldots\)。作为额外几何复核，曲线左侧真实面积 \(\int_0^2(1+y^2/4)dy=8/3\)，与题设经验值恰好一致。

\gkanswer{答案（1）：\(y^2=4x\)，\(0\le y\le2\)。}

\gkanswer{答案（2）：矩形面积 \(\dfrac52\)，五边形面积 \(\dfrac{11}{4}\)，五边形更接近 \(\dfrac83\)。}

\end{answer}
\begin{answer}
\gksolhead{21}{}

\textbf{题目}：

双曲线 \(x^2-\dfrac{y^2}{b^2}=1\)（\(b>0\)）的左、右焦点分别为 \(F_1\)、\(F_2\)，直线 \(l\) 过 \(F_2\) 且与双曲线交于 \(A\)、\(B\) 两点。

（1） 若 \(l\) 的倾斜角为 \(\dfrac\pi2\)，\(\triangle F_1AB\) 是等边三角形，求双曲线的渐近线方程。

（2） 设 \(b=\sqrt3\)，若 \(l\) 的斜率存在，且 \(|AB|=4\)，求 \(l\) 的斜率。

\Needspace{16\baselineskip}

\Needspace{4\baselineskip}
\textbf{解析}

\gksubpart{(1)}

答案（已纠正）：渐近线 \(y=\pm\sqrt2\,x\)。
\Needspace{5\baselineskip}
焦点 \(F_1=(-c,0),F_2=(c,0)\)，\(c^2=1+b^2\)。倾角 \(\pi/2\) 的直线为 \(x=c\)，代入原方程 \(x^2-y^2/b^2=1\)：
\nopagebreak[4]
\[
\frac{y^2}{b^2}=c^2-1=b^2,\qquad y^2=b^4.
\]
因 \(b>0\)，两交点为 \(A=(c,b^2),B=(c,-b^2)\)（标号可交换），不是 \((c,\pm bc)\)。
\Needspace{5\baselineskip}
故 \(|AB|=2b^2\)，\(|F_1A|=|F_1B|=\sqrt{4c^2+b^4}\)。等边条件为
\nopagebreak[4]
\[
4b^4=4c^2+b^4\iff3b^4=4(1+b^2).
\]
令 \(u=b^2>0\)，则 \(3u^2-4u-4=(3u+2)(u-2)=0\)，排除 \(u=-2/3\)，得 \(b^2=2\)、\(b=\sqrt2\)。因此渐近线为 \(y=\pm bx=\pm\sqrt2\,x\)。
精确回代：\(c=\sqrt3\)，\(A,B=(\sqrt3,\pm2)\)，原方程左端为 \(3-4/2=1\)；三边长度分别为 \(4,\sqrt{12+4},\sqrt{12+4}\)，均为4。置信度 H。

\gksubpart{(2)}

答案：\(k=\pm\sqrt{15}/5\)。此时 \(b=\sqrt3,c=2\)，直线写为 \(x=2+t,y=kt\)。代入 \(x^2-y^2/3=1\) 得
\((3-k^2)t^2+12t+9=0\)。
若 \(k^2=3\) 退化为一次方程，只有一个交点，不能满足题意。否则两根差的绝对值为 \(6\sqrt{1+k^2}/|3-k^2|\)，乘方向向量长度 \(\sqrt{1+k^2}\) 得
\(|AB|=6(1+k^2)/|3-k^2|=4\)。
令 \(u=k^2\ge0\)：当 \(u<3\)，\(6+6u=12-4u\)，得 \(u=3/5\)；当 \(u>3\)，\(6+6u=4u-12\) 无非负解。两种符号均有效。复算 \(6(1+0.6)/(3-0.6)=4\)，斜率绝对值约 \(0.7745966692\)。

\gkanswer{答案（1）：\(y=\pm\sqrt2\,x\)。}

\gkanswer{答案（2）：\(k=\pm\sqrt{\dfrac35}\)。}

\end{answer}
\begin{answer}
\gksolhead{22}{}

\textbf{题目}：

对于无穷数列 \(\{a_n\}\) 与 \(\{b_n\}\)，记 \(A=\{x\mid x=a_n,\ n\in\mathbf N^*\}\)，\(B=\{x\mid x=b_n,\ n\in\mathbf N^*\}\)，若同时满足条件：① \(\{a_n\}\)，\(\{b_n\}\) 均单调递增；② \(A\cap B=\varnothing\) 且 \(A\cup B=\mathbf N^*\)，则称 \(\{a_n\}\) 与 \(\{b_n\}\) 是无穷互补数列。

（1） 若 \(a_n=2n-1\)，\(b_n=4n-2\)，判断 \(\{a_n\}\) 与 \(\{b_n\}\) 是否为无穷互补数列，并说明理由；

（2） 若 \(a_n=2^n\) 且 \(\{a_n\}\) 与 \(\{b_n\}\) 是无穷互补数列，求数列 \(\{b_n\}\) 的前 16 项的和；

（3） 若 \(\{a_n\}\) 与 \(\{b_n\}\) 是无穷互补数列，\(\{a_n\}\) 为等差数列且 \(a_{16}=36\)，求 \(\{a_n\}\) 与 \(\{b_n\}\) 的通项公式。

\Needspace{4\baselineskip}
\textbf{解析}

无穷互补要求两列严格递增，且值集互不相交、合为正整数。

\gksubpart{(1)}

不是。\(a_n=2n-1\) 是正奇数，\(b_n=4n-2\) 是模4余2的正整数，交集为空，但4既不在A也不在B，故不覆盖正整数。

\gksubpart{(2)}

答案：前16项和为180。\(a_n=2^n\)（\(n\ge1\)），所以需删除 \(2,4,8,16,\ldots\)，保留1。\(b_1\) 至 \(b_{16}\) 依次为
\(1,3,5,6,7,9,10,11,12,13,14,15,17,18,19,20\)。
在1到20中删除 \(2,4,8,16\) 后恰余16个数，因此和为 \(20\cdot21/2-(2+4+8+16)=210-30=180\)。程序独立枚举复算同为180。

\gksubpart{(3)}

答案：\(a_n=2n+4\)；\(b_n=n\)（\(1\le n\le5\)），\(b_n=2n-5\)（\(n\ge6\)）。
设等差列 \(a_n=a_1+(n-1)d\)。严格递增且值均为正整数，故 \(a_1,d\) 都是正整数。由 \(a_{16}=36\) 得 \(a_1+15d=36\)，正整数解仅 \((a_1,d)=(21,1),(6,2)\)。第一种A含全部大于等于21的整数，补集只有1到20，不能组成无穷严格递增的B，排除。第二种A是所有大于等于6的偶数；补集是1到5及从7起的奇数，按递增排列即上述公式。核验 \(a_{16}=36,b_5=5,b_6=7\)，两列确实恰分正整数。

\gkanswer{答案（1）：不是无穷互补数列。}

\gkanswer{答案（2）：\(180\)。}

\gkanswer{答案（3）：\(a_n=2n+4\)；\(b_n=n\)（\(1\le n\le5\)）；\(b_n=2n-5\)（\(n\ge6\)）。}

\end{answer}
\begin{answer}
\gksolhead{23}{}

\textbf{题目}：

已知 \(a\in\mathbf R\)，函数 \(f(x)=\log_2\left(\dfrac1x+a\right)\)。

（1） 当 \(a=1\) 时，解不等式 \(f(x)>1\)；

（2） 若关于 \(x\) 的方程 \(f(x)+\log_2(x^2)=0\) 的解集中恰有一个元素，求 \(a\) 的值；

（3） 设 \(a>0\)，若对任意 \(t\in\left[\dfrac12,1\right]\)，函数 \(f(x)\) 在区间 \([t,t+1]\) 上的最大值与最小值的差不超过 1，求 \(a\) 的取值范围。

\Needspace{4\baselineskip}
\textbf{解析}

函数 \(f(x)=\log_2(1/x+a)\)，始终要求 \(x\ne0\) 且 \(1/x+a>0\)。

\gksubpart{(1)}

答案：\((0,1)\)。当 \(a=1\)，\(f(x)>1\) 等价 \(1/x+1>2\)，即 \(1/x>1\)。负数x不能满足；正数时得 \(x<1\)，故 \(0<x<1\)，该区间自动满足定义域。

\gksubpart{(2)}

答案：\(a\in\{-1/4,0\}\)。原方程在定义域内等价
\(\log_2[(1/x+a)x^2]=0\)，即 \(ax^2+x-1=0\)。
任何此方程的实根都不为0，且 \(1/x+a=1/x^2>0\)，故所有实根均回到原方程，没有增根。
当 \(a=0\) 时是一次方程，唯一根1；当 \(a\ne0\) 时恰有一个不同实根当且仅当判别式 \(1+4a=0\)，得 \(a=-1/4\)，重根2。判别式正时两根不同且均合法，负时无实根。
复核：\(a=0,x=1\) 两个对数均0；\(a=-1/4,x=2\)，\(f(2)=-2\)，\(\log_2(4)=2\)。

\gksubpart{(3)}

答案：\([2/3,+\infty)\)。由于 \(a>0\) 且 \(t\in[1/2,1]\)，区间 \([t,t+1]\) 全在正半轴，函数严格递减。最大最小差为
\(D(t)=\log_2\frac{a+1/t}{a+1/(t+1)}\)。
\(D(t)\le1\) 等价 \(a\ge1/t-2/(t+1)=(1-t)/[t(t+1)]\)。记右端为q(t)，其导数
\(q'(t)=(t^2-2t-1)/[t^2(t+1)^2]<0\)（\(1/2\le t\le1\)），故最大值 \(q(1/2)=2/3\)。因此对全部t成立恰需 \(a\ge2/3\)。边界复算 \(a=2/3,t=1/2\) 时比值 \((2/3+2)/(2/3+2/3)=2\)，差恰为1。

\gkanswer{答案（1）：\((0,1)\)。}

\gkanswer{答案（2）：\(a\in\{-\dfrac14,0\}\)。}

\gkanswer{答案（3）：\([\dfrac23,+\infty)\)。}

\end{answer}